\par\bigskip
\SourceLeaf{421}
\begin{center}
\pdfbookmark[2]{Gauss an Olbers, 26. November 1810}{pallas-brief-18101126}
\textsc{Gauss} an \textsc{Olbers}, 26. November 1810.
\end{center}
\noindent\hspace*{\parindent}Ich glaube, ich habe Ihnen bereits in meinem letzten Brief geschrieben, dass ich eine Arbeit\HistoricalNoteEight{Vgl. S. 473 dieses Bandes.]} über die Störungen der Pallas durch Jupiter unternehmen wollte. Dabei wollte ich allein ihren Betrag von 1803 bis 1811 berechnen und ihre allgemeine Theorie beiseitelassen. Die Berechnung dreier Sätze elliptischer Elemente aus den Oppositionen von 1803, 1804, 1805, 1807 — 1804, 1805, 1807, 1808 — 1805, 1807, 1808, 1809 hatte die starke Wirkung der Störungen gezeigt. Urteilen Sie selbst anhand der folgenden Beispiele:
\begin{center}
\begingroup\footnotesize\setlength{\tabcolsep}{3pt}\renewcommand{\arraystretch}{1.2}
\begin{tabular}{lccc}
Perihel 1803\ .\,.\,.\,.\,.\,.\,.\,. & \(121^{0}\,3'11''\) & \(121^{0}\,5'22''\) & \(120^{0}58'\,5''\)\\
\(\varphi\)\ .\,.\,.\,.\,.\,.\,.\,.\,.\,.\,.\,.\,.\,.\,.\,.\,.\,.\,.\,.\,. & \(14^{0}10'59''\) & \(14^{0}10'\,4''\) & \(14^{0}\,9'37''\)\\
Tägl. mittl. Bewegung.\ .\,.\,. & \(770''{,}214\) & \(770''{,}447\) & \(770''{,}926.\)\\
\end{tabular}\endgroup
\end{center}
\par
\noindent\hspace*{\parindent}Ich wurde neugierig, ob sich diese sechs Oppositionen viel besser miteinander vereinbaren ließen, wenn man die Störungen durch Jupiter berücksichtigte. Ich gestehe, dass eine solche an sich unangenehme Arbeit für mich zu wenig Reiz hatte, solange eine einzige Ellipse alle Beobachtungen darstellte, und dass mich selbst der Preis von 6000 Francs kaum dazu hätte verleiten können. Jetzt habe ich die Arbeit jedoch mutig in Angriff genommen. Nachdem ich ihr 4 Wochen lang den größten Teil meiner verfügbaren Zeit gewidmet habe, ist die Hauptsache erledigt, sodass ich, wenn meine Rechnung keine wesentlichen Fehler enthält, bald die Früchte davon zu ernten hoffe. Ich hielt es für am zweckmäßigsten, die Elemente als veränderlich zu betrachten und ihre täglichen Änderungen durch Jupiter von 50 zu 50 Tagen über einen Zeitraum von mehr als 8 Jahren zu berechnen. Ich begann damit, eigene Formeln dafür zu entwickeln, oder vielmehr die bereits vor mehreren Jahren von mir entwickelten in eine möglichst handliche Form zu bringen. Die Formeln von \textsc{Laplace} erschienen mir viel zu umständlich, und bei ihrer Anwendung hätte ich gewiss zwei- bis dreimal so viel Zeit gebraucht wie mit meinen. Die Positionen des Jupiters berechnete ich nach den Tafeln von \textsc{Bouvard} (mit einem besonderen Kunstgriff\HistoricalNoteDouble{Vgl. S. 475 dieses Bandes.]}) \SourceLeaf{422}und die Positionen der Pallas nach einem 4. Satz von Elementen, der sich möglichst genau an alle Oppositionen anschließt. Heute habe ich endlich jene 61 täglichen Änderungen für jedes der sechs Elemente fertig berechnet. Übrig bleiben die mechanische Integration (die kaum mehr als Addieren ist) und die Neuberechnung der Grundelemente, die zusammen mit den sich so ergebenden Störungen die Oppositionen am besten darstellen. Diese Störungen sind viel beträchtlicher, als ich dachte, und besonders stark im Jahr 1809, als Jupiter und Pallas in heliozentrischer Konjunktion standen. Ich glaube deshalb, dass meine Ephemeride bei der nächsten Opposition sehr stark abweichen wird. \textsc{Harding} hat neulich nach Pallas gesucht, weiß aber noch nicht sicher, ob er sie gefunden hat. Sollten Sie sie beobachten, wäre ich Ihnen für eine baldige Mitteilung sehr verbunden. Die tägliche Bewegung wird jetzt um mehr als \(2''\) größer sein als 1803. Der Winkel \(\varphi\), der beständig abgenommen hat, wird wohl um \(13'\) kleiner sein, und das Perihel ist wohl auch um \(20'\) oder \(30'\) zurückgegangen usw. Sollte die Übereinstimmung mit den sechs bisherigen Oppositionen und der 7. des nächsten Jahres zu meiner Zufriedenheit ausfallen, nehme ich vielleicht auch die allgemeine Theorie in Angriff, obwohl mich das ungeheure Ausmaß an mechanischer Rechnung dabei sehr abschreckt. Die Störungen durch Jupiter überwiegen so sehr, dass man die durch die anderen Planeten vorerst beiseitelassen kann.\quad\mbox{.\quad.\quad.\quad.\quad.}
\par
\noindent\hspace*{\parindent}Sollten Sie durch einen Hinweis in Paris eine weitere Verlängerung der Frist für die Preisfrage bewirken, wäre das für mich ein Grund mehr, die eben erwähnte Arbeit in Angriff zu nehmen. Zuerst muss mir jedoch die gegenwärtige gelingen. In vierzehn Tagen werde ich darüber im Klaren sein, zumindest in Bezug auf die bisherigen Beobachtungen. Doppelt neugierig werde ich dann auf die bevorstehenden sein, an denen sich die bei der Konjunktion erlittenen Störungen erst in ihrer ganzen Stärke zeigen werden.\quad.\quad.\quad.\quad.\quad.
\par
\Needspace{15\baselineskip}
\noindent\hspace*{\parindent}P. S. Vor dem Schluss dieses Briefs habe ich noch die Integration der Störung der Neigung abgeschlossen, die am wenigsten beträchtlich ist. Setze ich die Neigung bei der 1. Opposition
\begin{center}1803\quad auf\quad \(34^{0}37'47'',\)\end{center}
so beträgt sie
\SourceLeaf{423}
\begin{center}
\begin{tabular}{lc}
1804 & \(34^{0}37'36''\)\\
1805 & \(34^{0}37'12''\)\\
1807 & \(34^{0}37'25''\)\\
1808 & \(34^{0}37'31''\)\\
1809 & \(34^{0}36'40''\)\\
1811 & \(34^{0}35'\,5''.\)\\
\end{tabular}
\end{center}
\par
